\chapter{机器如何学习}
% 完整版：chapters/06_dofa.tex

到目前为止，所有电路里连接的强度都是工程师调好的。活的大脑里没有工程师。连接必须在工作过程中自己调整——而且要调得让机器做点有用的事。这就是学习。

最主要的限制听起来非常严格。每个连接只能自己改变自己，而它能知道的只有身边发生的事：发送者是否在工作，接收者是否在工作，以及有哪些物质到达了它。谁也不能给某个连接单独寄去一份修正。这就排除了训练人工神经网络的主要方法：在那里，误差沿网络反向传播，每条连接都得到专属于自己的修正。

\section*{一起工作，就连在一起}

最简单的规则是：如果发送者和接收者一起工作，它们之间的连接就会增强。这条规则是局部的，也不需要任何全局同步。只要再加上一点遗忘，这样的规则就能从信号流中找出最主要的东西——变化最大的那部分。还有一个看咔嗒时刻的版本：如果发送者在接收者\emph{之前}咔嗒，也就是说它可能是原因，连接就增强；如果在之后，连接就减弱。

但这些规则有一个局限：它们学到的是\emph{常见的}，而不是\emph{有用的}。一个只看得见邻居的连接，并不知道这个动作带来的是果汁还是疼痛。

\section*{第三个信号}

它由多巴胺带来——那个播报“比预期的好”的电台。让连接只在三件事同时发生时才改变：发送者工作了，接收者工作了，多巴胺来了。

\pencilfig{6}{\pencilart{6}}{当两个神经元一起工作过、并且第三个信号到来——“结果比预期的好”——时，连接就会增强。}

但有个难处：奖励不是马上就来，而是要过一两秒。到那时，发送者和接收者早就沉默了。连接需要记住自己曾经参与过。它确实记得：一起工作会在连接上留下一个标记，像一张一两秒后自己就会脱落的便利贴。便利贴还贴着的时候多巴胺来了——连接就改变。没来——标记便消失得无影无踪。这解决了寻址问题：多巴胺人人一份，但只有贴着便利贴的连接才会改变。

\pencilfig{106}{%
\foreach \y/\d/\t in {0/0.9/{及时：连接增强},-1.3/3.3/{来晚了：无变化}}{
  \draw[pen] (0,\y) -- (4.6,\y);
  \foreach \s in {0.25,0.33}{\draw[pcBlue, line width=0.8pt] (\s,\y) -- (\s,\y+0.3);}
  \fill[pcAmber, opacity=0.25] (0.35,\y) -- plot[domain=0.35:4.5, samples=30] (\x,{\y+0.8*exp(-(\x-0.35)/0.8)}) -- (4.5,\y) -- cycle;
  \draw[penlite=pcAmber] plot[domain=0.35:4.5, samples=30] (\x,{\y+0.8*exp(-(\x-0.35)/0.8)});
  \draw[pcAmber!80!black, line width=0.9pt] (\d-0.1,\y+0.95) -- (\d+0.07,\y+0.62) -- (\d-0.07,\y+0.6) -- (\d+0.1,\y+0.22);
  \draw[arr=pcAmber!80!black] (\d+0.09,\y+0.27) -- (\d+0.11,\y+0.12);
  \node[lbl, anchor=west] at (4.7,\y+0.3) {\t};}
\node[lbl, text=pcBlue, anchor=south west] at (0.1,0.85) {一起工作过};
\node[lbl, text=pcAmber!80!black, anchor=south] at (2.3,0.35) {标记消退};
\foreach \x/\l in {0.35/0,1.95/1,3.55/2}{\draw[pcGray, line width=0.5pt] (\x,-1.3) -- (\x,-1.38); \node[lbl, anchor=north] at (\x,-1.38) {\l~s};}
}{一起工作会在连接上留下一个标记，一两秒内消退。多巴胺只改变带有标记的连接。}

\section*{把噪声当作搜索}

为什么这条规则会让事情变好，而不是乱变？想象你蒙着眼睛寻找山顶。你随机迈出一步，有人喊“高了！”或“低了！”。如果是“高了”——就把这一步固定下来。一步又一步，你就爬了上去，一眼地图都没看。大脑做的正是同一件事：神经元带着一点随机的抖动，多巴胺报告情况是否变好，参与过的连接就朝有利的方向挪一点。原本妨碍传递消息的噪声，在这里变成了搜索。

现在就明白了，为什么多巴胺报告的不是奖励，而是\emph{意外}。如果它只是在每次成功后都喊“果汁！”，它会把网络当时做的一切都一并强化——对的和错的都一样。在我们的模型里，这样的网络确实会把连接放大几千倍，有时还会永远卡在最差的选择上。减去预期，才让学习变得公正。

\section*{一根导线的代价}

广播式的学习是有代价的。所有人共用一个信号，而信号里混着所有影响结果的神经元的噪声。网络越大，每个连接要分辨出属于自己的那一份就越久。因此，多巴胺看来只训练比较小的、负责选择动作的回路，而庞大的皮层主要依靠局部规则学习。

\section*{谁来计算意外}

还剩一个问题：多巴胺神经元怎样算出“比预期的好”？需要一个评论家——一群神经元，时时刻刻估计前面还有多少好事。意外等于得到的奖励加上预期跳升了多少。灯亮了——预期跳了上去，爆发。预告过的果汁来了——奖励是有了，但预期在同一刻被用完，没有意外。果汁没来——预期白白落空，低谷。第一章里的三种情况全都自然而然地出现了。多巴胺确实这样表现，这是测量到的。大脑究竟如何算出这个意外——目前还是假说。

\fullversion{6}
